\documentclass[sigconf, nonacm]{acmart}

\newcommand\vldbdoi{XX.XX/XXX.XX}
\newcommand\vldbpages{XXX-XXX}
\newcommand\vldbvolume{20}
\newcommand\vldbissue{1}
\newcommand\vldbyear{2026}
\newcommand\vldbauthors{\authors}
\newcommand\vldbtitle{\shorttitle}
\newcommand\vldbavailabilityurl{https://github.com/emailvenkatm/kndb/tree/postgres-experiment}
\newcommand\vldbpagestyle{plain}

\usepackage{booktabs}
\usepackage{listings}
\usepackage{microtype}
% Long inline \texttt{} identifiers otherwise trip overfull \hbox warnings;
% these settings let TeX break inside typewriter tokens and stretch spacing.
\sloppy
\hbadness=5000
\hfuzz=6pt
\lstset{
  basicstyle=\ttfamily\footnotesize,
  breaklines=true,
  columns=fullflexible,
  keepspaces=true,
  frame=single,
  framesep=3pt
}

\begin{document}

\title{Epistemic Typing as a PostgreSQL Table Access Method: Adversarial Conflict Resolution Under Confidence Forgery and Sybil Coordination}

\author{[Author Name Placeholder]}
\affiliation{%
  \institution{[Institution Placeholder]}
  \city{[City Placeholder]}
  \country{[Country Placeholder]}
}
\email{[email@placeholder.example]}

\begin{abstract}
We describe KNDB, a PostgreSQL 18 table access method (TAM) that types every row with an engine-assigned epistemic kind (MEASURED, INFERRED, or DERIVED) and resolves per-slot conflicts inside every write-time \texttt{heapam} callback. Rows land as ordinary heap tuples; seven of the 42 TAM callbacks are overridden (\texttt{tuple\_insert}, \texttt{multi\_insert}, \texttt{tuple\_update}, \texttt{tuple\_delete}, \texttt{tuple\_insert\_speculative}, \texttt{tuple\_complete\_speculative}, \texttt{relation\_toast\_am}), the other 35 delegate to heap; we provide a completeness argument over the interface as a paper artefact. This paper reports the engineering behind that decision and the adversarial evaluation that motivated it. On a confidence-forgery workload where an attacker asserts INFERRED writes with confidence in $[0.95, 1.0]$ against honest MEASURED writes with confidence in $[0.5, 0.9]$, KNDB beats a confidence-only baseline by 63 percentage points on the Book-Author fusion dataset and 92.7 points on the Zheng crowdsourcing dataset. Both wins are proven load-bearing on the kind axis by a source-rebuild disable-and-test in which the lattice is neutralised and the win vanishes. Against four truth-discovery baselines (TruthFinder, CRH, CATD, ACCU) reimplemented from the original equations and validated to within 0.3 percentage points of the published numbers, KNDB is competitive below a per-dataset density-saturation cell and dominant at or above it. We formalise the cell as $k^{*} \approx \rho_{\text{alg}} \cdot h_{\text{top}}$, where $h_{\text{top}}$ is per-slot top honest surface-form support, and validate the prediction within $\pm 20\%$ on Book-Author and $\pm 30\%$ on Zheng. Because the kind axis is assigned by the engine from independent metadata and cannot be forged at write time, KNDB's $k^{*}$ is unbounded. The paper is honest about where KNDB loses: CRH and ACCU outperform KNDB below saturation on Zheng, and KNDB scores zero on three temporal knowledge-editing benchmarks whose ground truth is last-writer-wins.
\end{abstract}

\maketitle

\pagestyle{\vldbpagestyle}

%%% ---------------------------------------------------------------
\section{Introduction}
%%% ---------------------------------------------------------------

Databases now ingest writes from three broad classes of producer that disagree systematically about what they know. Instruments and user-facing forms produce claims that are directly observed. Machine-learning models and LLM agents produce claims that are inferred from a prompt or a feature vector. Rule engines produce claims that are derived from other rows already in the database. A conventional relational store treats these writes identically: whatever arrives last, or whatever a user-space \texttt{BEFORE INSERT} trigger accepts, becomes the current value. Under adversarial pressure this is unsafe. A miscalibrated LLM writer can flood a slot with high-confidence hallucinations that a confidence-ranking arbitrator will prefer over honest instrument readings. A coordinated set of Sybil agents can outvote honest sources on any slot where the Sybil count exceeds the top honest surface-form support.

KNDB is a PostgreSQL 18 extension that responds by lifting the epistemic kind of each write into the storage engine itself. Every row carries an \texttt{ep\_kind} column (MEASURED, INFERRED, or DERIVED); at \texttt{tuple\_insert} time a table access method (TAM) hook applies a total order on (kind, specificity, confidence, xmin) and evicts any live incumbent that loses. Because the hook runs inside \texttt{heapam}, no user-space bypass (\texttt{DISABLE TRIGGER}, \texttt{session\_replication\_role = 'replica'}, \texttt{COPY FROM}) reaches around it. The engine is the enforcement point, not the trigger.

The central quantitative claim is Sybil-robustness at 1:1 density. Every truth-discovery (TD) baseline we tested collapses when the Sybil count $k$ reaches or exceeds the per-slot top honest surface-form support $h_{\text{top}}$. KNDB does not, because kind rank is assigned by the engine from independent metadata that the adversarial identity does not possess (a source's \texttt{n\_listings} history on Book-Author, a worker's \texttt{quali\_acc} on a disjoint qualification test on Zheng). We prove this both empirically (Section~\ref{sec:evaluation}) and as a threshold theorem (Section~\ref{sec:formalism}).

We are careful about scope. On workloads where kind and confidence coincide by construction, the lattice's kind axis buys nothing beyond confidence sorting, and KNDB matches \texttt{pg\_conf} exactly (Section~\ref{sec:evaluation}). On temporal knowledge-editing benchmarks whose ground truth is last-writer-wins, KNDB scores zero at $c=1$, because its \texttt{xmin} tiebreak is first-committer-wins by design (Section~\ref{sec:limitations}). Below the density-saturation cell, CRH and ACCU outperform KNDB on the Zheng workload. These are stated up front and detailed in Sections~\ref{sec:evaluation}~and~\ref{sec:limitations}.

\textbf{Contributions.}
\begin{enumerate}
\item An engine-level epistemic type system implemented as a PostgreSQL 18 TAM. Three callbacks overridden (\texttt{tuple\_insert}, \texttt{multi\_insert}, \texttt{relation\_toast\_am}); the rest delegate to heap. Total 1789 lines: 1479 of C in \texttt{src/} and 310 of headers in \texttt{include/} (\texttt{wc -l src/*.c include/*.h}).
\item An enumerated bypass model. We name six user-space write paths (\texttt{DISABLE TRIGGER}, \texttt{session\_replication\_role}, \texttt{COPY FROM} single-row and batch, \texttt{INSERT}, apply worker) and one that is out of scope (\texttt{ALTER TABLE \dots SET ACCESS METHOD heap}). All six write paths are blocked by the TAM; the schema-change threat is disclosed. The batch \texttt{COPY} path required a dedicated \texttt{multi\_insert} override to close, described in Section~\ref{sec:implementation}.
\item A confidence-forgery workload on two independent datasets, with source-rebuild disable-and-test in each case. Every quantitative claim in this paper corresponds to a JSON in \texttt{bench/results/} whose dylib SHA-256 is recorded in the run metadata.
\item A threshold theorem $k^{*} \approx \rho_{\text{alg}} \cdot h_{\text{top}}$ that predicts the Sybil-collapse cell of any dataset from its per-slot honest-support distribution alone. Empirically validated within $\pm 20\%$ on Book-Author (predicted bracket [3,5], observed cliff N=5\ldots10) and $\pm 30\%$ on Zheng (predicted bracket [14,20], observed cell N=20).
\item Honest limitations. TD baselines legitimately beat KNDB below saturation; KNDB scores 0.000 on three temporal benchmarks; there is a 15\,000-row-per-transaction ceiling at PostgreSQL's default \texttt{max\_locks\_per\_transaction=64}, inherited into \texttt{COPY}; \texttt{ALTER TABLE \dots SET ACCESS METHOD heap} sits outside the enforcement envelope; the custom WAL rmgr is an annotation channel, not a durability channel.
\end{enumerate}

Source, benchmark harnesses, and every result JSON cited below are available at \url{https://github.com/emailvenkatm/kndb/tree/postgres-experiment} (branch \texttt{postgres-experiment}, tip \texttt{22d98ed} at time of writing). Each result JSON in \texttt{bench/results/} carries hardware and PostgreSQL configuration metadata, and each cell that involved a source rebuild records the dylib SHA-256 in its metadata block.

%%% ---------------------------------------------------------------
\section{Motivation and threat model}
\label{sec:threat}
%%% ---------------------------------------------------------------

The threat model is stated up front and does not adapt to outcomes.

\subsection{Confidence forgery}
The confidence-forgery adversary asserts INFERRED writes with self-reported confidence uniform in $[0.95, 1.0]$ against a slot whose lattice-correct incumbent is a Tier-A MEASURED writer with self-reported confidence uniform in $[0.5, 0.9]$. Adversarial payload contains a scrambled real answer (an author string lifted from a different gold ISBN in Book-Author; a flipped binary label in Zheng). The real-world analogue is an LLM-generated claim that hallucinates a value but self-reports as certain. A confidence-only arbitrator ranks by confidence and prefers the adversary's higher self-report. KNDB is designed to reject the adversary regardless of the self-reported confidence, because MEASURED strictly outranks INFERRED in the lattice.

\subsection{Sybil coordination}
The Sybil adversary controls $k$ synthetic identities per gold slot, each asserting the same falsified value. The adversary knows gold and picks the falsified value to maximise confusion (in binary tasks this is uniquely determined). The adversary cannot forge existing honest identities (Zheng worker IDs, Book-Author bookstore names) and cannot forge a MEASURED kind assignment. Kind assignment is engine-controlled and requires metadata that adversarial identities have no history of producing: on Book-Author, tier rank uses \texttt{n\_listings} (source volume in \texttt{book.txt}) and \texttt{canon\_rate} (fraction of author fields in "Last, First" format); on Zheng, tier rank uses \texttt{quali\_acc}, computed from a disjoint qualification test with question IDs 2000\ldots 2019 that every worker took. Both signals live in the dataset before the conflict resolution runs, and both are unavailable to an adversary that appeared for the first time in the write stream. The independence self-audits are in the datasets' READMEs.

\subsection{Write-path bypass}
The engine-in-storage claim rides on the TAM callback being reachable from every user-space write path that a hostile writer with \texttt{INSERT} plus \texttt{ALTER TABLE} on the target relation can construct. Enumerated write paths, each verified against PostgreSQL 18 REL\_18\_STABLE source:
\begin{itemize}
\item \texttt{ALTER TABLE \dots DISABLE TRIGGER ALL} sets \texttt{pg\_trigger.tgenabled = 'D'} (\texttt{tablecmds.c:5588-5592}); \texttt{TriggerEnabled} at \texttt{trigger.c:3491-3499} returns false. A trigger-based enforcer is bypassable here; a TAM callback is not. Note that \texttt{TriggerEnabled} at \texttt{trigger.c:3489-3499} still skips triggers marked \texttt{TRIGGER\_DISABLED} even when \texttt{session\_replication\_role = 'replica'}, so \texttt{DISABLE TRIGGER ALL} defeats \texttt{ENABLE ALWAYS TRIGGER} as well; the AM callback fires regardless of any \texttt{pg\_trigger.tgenabled} value.
\item \texttt{SET session\_replication\_role = 'replica'}. \texttt{TriggerEnabled} at \texttt{trigger.c:3489-3499} skips \texttt{TRIGGER\_FIRES\_ON\_ORIGIN} (the default for user \texttt{CREATE TRIGGER}) under \texttt{SESSION\_REPLICATION\_ROLE\_REPLICA}. Triggers created \texttt{ENABLE ALWAYS} (\texttt{tablecmds.c:5566}, \texttt{TRIGGER\_FIRES\_ALWAYS}) are not in the disable-set at either branch of \texttt{TriggerEnabled} and DO fire under \texttt{replica}; that closes this specific bypass for a trigger-based enforcer that opts in to \texttt{ENABLE ALWAYS}. The GUC is \texttt{PGC\_SUSET} (\texttt{guc\_tables.c:5166}). The trigger-based enforcer therefore requires \texttt{ENABLE ALWAYS} to survive the \texttt{replica} bypass and still cannot survive \texttt{DISABLE TRIGGER ALL}; the AM callback survives both.
\item \texttt{COPY FROM} single-row path. \texttt{CopyFrom} at \texttt{copyfrom.c:995-1006} forces \texttt{CIM\_SINGLE} when the target has a BEFORE trigger, an FDW that does not batch, or partitioned tables with statement-level triggers. The single-row path at \texttt{copyfrom.c:1427} calls \texttt{table\_tuple\_insert}, which routes through the TAM.
\item \texttt{COPY FROM} batch path (\texttt{CIM\_MULTI}). The default path with no BEFORE trigger. Batched rows flow through \texttt{table\_multi\_insert} at \texttt{copyfrom.c:554-559}, dispatching on the AM's \texttt{multi\_insert}. Before we overrode it (Section~\ref{sec:implementation}), this inherited \texttt{heap\_multi\_insert} and silently skipped every epistemic rule check.
\item \texttt{INSERT}, \texttt{INSERT \dots SELECT}, \texttt{INSERT \dots VALUES}. Standard \texttt{ModifyTable} path lands in \texttt{ExecInsert}, which calls \texttt{table\_tuple\_insert}.
\item \texttt{UPDATE} that rewrites the epistemic prefix (F20). \texttt{ModifyTable} on the \texttt{CMD\_UPDATE} branch reaches \texttt{table\_tuple\_update}, which dispatches on the AM's \texttt{tuple\_update} callback (\texttt{tableam.h:718-727}). Before F20 we inherited \texttt{heapam\_tuple\_update} (\texttt{heapam\_handler.c:392-400}, delegating to \texttt{heap\_update} at \texttt{heapam.c:3241}) and an \texttt{UPDATE fact SET ep\_kind='MEASURED', ep\_confidence=1.0} against an incumbent \texttt{INFERRED/0.4} row committed cleanly, forging the epistemic prefix on a live row without R1--R5 firing, without the advisory lock, and without the precedence lattice or eviction audit. The F20 override refuses any \texttt{UPDATE} whose new prefix differs from the incumbent's; content-only updates (value, valid\_time, sources) are permitted.
\item \texttt{DELETE} that removes an incumbent (F20). \texttt{ModifyTable} on \texttt{CMD\_DELETE} reaches \texttt{table\_tuple\_delete}, which dispatches on the AM's \texttt{tuple\_delete} callback (\texttt{tableam.h:709-716}). Before F20 we inherited \texttt{heapam\_tuple\_delete} (\texttt{heapam\_handler.c:382-388}, delegating to \texttt{heap\_delete} at \texttt{heapam.c:2772}). An adversary could \texttt{DELETE} a \texttt{MEASURED} incumbent and then \texttt{INSERT} an \texttt{INFERRED} forgery unopposed, since the precedence lattice has nothing left to compare against; no eviction audit row would ever be written for the \texttt{MEASURED} loss. The F20 override refuses \texttt{DELETE} outright with \texttt{ERRCODE\_FEATURE\_NOT\_SUPPORTED}.
\item \texttt{INSERT \dots ON CONFLICT} (speculative-insertion protocol, F21). \texttt{ExecInsert} at \texttt{nodeModifyTable.c:1189-1216} REL\_18\_STABLE handles \texttt{ON CONFLICT} by acquiring a \texttt{SpeculativeInsertionLockAcquire}, calling \texttt{table\_tuple\_insert\_speculative} (\texttt{tableam.h:513-519}) to write the tuple with a speculative token, probing arbiter indexes via \texttt{ExecInsertIndexTuples}, then calling \texttt{table\_tuple\_complete\_speculative} (\texttt{tableam.h:521-525}) to either confirm (\texttt{heap\_finish\_speculative}) or kill (\texttt{heap\_abort\_speculative}) the tuple. Before F21 the AM inherited \texttt{heapam\_tuple\_insert\_speculative} and \texttt{heapam\_tuple\_complete\_speculative} verbatim (\texttt{heapam\_handler.c:2639-2640}); an \texttt{INSERT \dots ON CONFLICT DO UPDATE} or \texttt{DO NOTHING} that reached the speculative path would bypass R1--R5, the advisory lock, precedence, and the eviction audit. \emph{The SQL surface for this attack is currently unreachable on epistemic tables} because \texttt{CREATE UNIQUE INDEX}, \texttt{ADD PRIMARY KEY}, \texttt{ADD UNIQUE}, and \texttt{ADD EXCLUDE} all fail at \texttt{heap\_getnext}'s \texttt{rd\_tableam} identity check (\texttt{heapam.c:1352}) during the arbiter index's build scan, so \texttt{ON CONFLICT (col)} has no arbiter to bind (see \texttt{scripts/speculative\_forgery.sh} for the empirical SQL-level probe). The F21 overrides are defensive: the engine-in-storage claim must not depend on that accidental index-support gap persisting. The load-bearing proof lives in a C-level probe (\texttt{epistemic.\_probe\_speculative\_insert}) and its \texttt{sql/am\_speculative.sql} regression cell, exercised by the source-rebuild disable-and-test in Section~\ref{sec:eval-f21}.
\item Logical replication apply worker. Runs as \texttt{SESSION\_REPLICATION\_ROLE\_REPLICA} by default; still reaches the TAM via \texttt{ExecSimpleRelationInsert} $\to$ \texttt{ExecInsert} $\to$ \texttt{table\_tuple\_insert}. Enforcement fires.
\item \texttt{ALTER TABLE \dots SET ACCESS METHOD heap}. Out of scope. This is a schema-change threat available to the table owner only; it rewrites the relation onto plain heap and removes the TAM from the write path entirely. See Section~\ref{sec:limitations}.
\end{itemize}

The engine-in-storage differentiator holds against a writer with \texttt{INSERT}, \texttt{UPDATE}, \texttt{DELETE}, plus \texttt{ALTER TABLE} on the target relation, or a role that can toggle \texttt{session\_replication\_role} (a replication or CDC operator, a migration tool, a connection pool). It does not hold against the table owner.

%%% ---------------------------------------------------------------
\section{Design}
\label{sec:design}
%%% ---------------------------------------------------------------

KNDB adds one PostgreSQL access method (\texttt{CREATE ACCESS METHOD epistemic USING TABLE HANDLER epistemic\_am\_handler}) and one base type (\texttt{epistemic\_kind}, a pass-by-value byte). Tables declared \texttt{USING epistemic} carry three extra columns beyond the user schema (\texttt{ep\_kind}, \texttt{ep\_specificity}, \texttt{ep\_confidence}) and a bitemporal \texttt{sys\_time tstzrange}. Rows are stored as ordinary heap tuples; storage format on disk is identical to a heap table with the same schema.

\subsection{Precedence lattice}
At write time a candidate row and any live incumbent for the same (entity\_id, attribute) slot are ranked by the total order:
\[
\underbrace{\text{kind}}_{\text{MEASURED}>\text{DERIVED}>\text{INFERRED}} \succ \text{specificity} \succ \text{confidence} \succ \text{xmin}
\]
Ranks are integer-compared at each level. Kind rank is the primary sort key. The lattice is defined in \texttt{epistemic\_rules.c} and exposed as an idempotent SQL function tested by \texttt{sql/precedence.sql}.

\subsection{Write-time rules R1--R5}
The TAM enforces five predicates on the candidate before ranking, checked in order at the top of \texttt{epistemic\_tuple\_insert\_impl}; the first failure ereports \texttt{ERRCODE\_CHECK\_VIOLATION} and aborts the insert. Reference implementation: \texttt{contrib/epistemic/src/epistemic\_rules.c}.
\begin{itemize}
\item R1 (DERIVED sources): if the candidate's kind is DERIVED, its \texttt{sources[]} column must be non-empty. A DERIVED assertion without a source is refused.
\item R2 (source resolution): if the candidate's kind is non-MEASURED and \texttt{sources[]} is non-empty, every source identifier must resolve against \texttt{epistemic.source\_registry}. R2 uses SPI to look up each source id and refuses the write if any is unregistered. This is the registry that prevents an adversary from inventing new source identities on the fly.
\item R3 (MEASURED no sources): if the candidate's kind is MEASURED, \texttt{sources[]} must be empty. A MEASURED row is a first-hand observation and does not derive from other sources.
\item R4 (INFERRED confidence bound): if the candidate's kind is INFERRED, \texttt{ep\_confidence} must be strictly less than 1.0. An INFERRED row that claims certainty is refused; certainty is the property MEASURED asserts.
\item R5 (slot kind registry): if \texttt{epistemic.slot\_kind} has a row for the candidate's \texttt{attribute}, the candidate's kind must match \texttt{required\_kind}. R5 lets an operator pin an attribute to a specific kind (e.g., ``a hemoglobin\_A1c reading must be MEASURED, never INFERRED'').
\end{itemize}
R1 through R5 are content-level constraints on individual candidate rows and are independent of the precedence lattice; a candidate that passes R1--R5 still competes against any live incumbent through the precedence check that follows. The confidence-forgery attack we exercise in Section~\ref{sec:eval} uses INFERRED writes with \texttt{ep\_confidence} in $[0.95, 1.0)$: adversarial writes are chosen precisely so they pass R4 (and R1, R2, R3, R5), forcing the kind axis of the precedence lattice to be the mechanism that defeats them.

\subsection{Per-slot advisory lock (F6)}
Between the R1--R5 check and the incumbent scan, the TAM takes a \texttt{LOCKTAG\_ADVISORY} transaction-scope lock with \texttt{(field2=entity\_id, field3=hash\_bytes(attribute), field4=2)}, constructed inline via \texttt{SET\_LOCKTAG\_ADVISORY} and acquired with \texttt{LockAcquire(\&tag, ExclusiveLock, false, false)}. This is the same tag, mode, and scope as \texttt{pg\_advisory\_xact\_lock\_int4}. It serialises concurrent writers on their shared (entity\_id, hashed attribute) key, so the incumbent scan runs against a snapshot that has already absorbed any peer's just-committed row. \texttt{find\_live\_overlap} and \texttt{epistemic\_close\_sys\_time} both use \texttt{GetLatestSnapshot()} for the same reason.

Without this lock, two concurrent same-slot writers each see an empty slot in their own MVCC snapshot, and both commit. The rc\_invariant.sh script runs 50 trials of two overlapping same-slot inserts. Under the honest build, session1 wins all 50 and no both-live rows land. With the advisory-lock block patched to \texttt{if (0)} and the dylib rebuilt, both writers commit in 50 of 50 trials.

\subsection{First-committer-wins tiebreak (F8)}
On a true (kind, specificity, confidence) tie the TAM decides by reading the incumbent's raw \texttt{xmin} via \texttt{HeapTupleHeaderGetRawXmin} and comparing against the current backend's xid via \texttt{TransactionIdPrecedes} (which handles xid wraparound). Under the advisory lock, the incumbent's transaction has already committed by the time the loser scans it, so \texttt{incumbent\_xmin < new\_xid} on every race. The incumbent wins.

The earlier design used a content-hash tiebreak. That gave up grind-resistance in exchange for content-determinism: an attacker with \texttt{SELECT} plus \texttt{INSERT} could iterate byte-level variations of \texttt{value} until \texttt{hash\_bytes} placed the attacker below the incumbent, and \texttt{scripts/hash\_grind.sh} at F7 reported 30 of 30 attacker wins with a mean of 7.5 attempts (min 1, max 95). Under the xmin tiebreak the same script reports 0 of 20\,000 attacker wins across 100 trials of 200 attempts. The tradeoff we accepted: survivor is start-order-dependent and is not stable across \texttt{pg\_dump}/\texttt{pg\_restore}, because restore reloads rows via \texttt{COPY FROM} at \texttt{copyfrom.c:1427} and each reloaded row gets a fresh xid. What is stable across dump and restore is the "exactly one live row per slot" invariant that \texttt{sql/am\_eviction.sql} asserts.

\subsection{Reason codes and audit}
On eviction the TAM writes one audit row into \texttt{epistemic.evicted\_fact} via SPI and closes the incumbent's \texttt{sys\_time} upper bound via \texttt{simple\_heap\_update}. The audit row carries the losing row's identifiers and a reason code from \texttt{EP\_REASON\_\{OUTRANKED\_KIND, OUTRANKED\_SPEC, OUTRANKED\_CONF, CONTRADICTED\_SAME\_RANK\}}. All three state changes (winner insert, audit insert, incumbent update) execute inside one top-level PostgreSQL transaction created by \texttt{start\_xact\_command}; the TAM does not open, commit, or manage transactions. Atomicity is inherited from PostgreSQL. An adversarial control that stages the audit row through \texttt{dblink} (which commits in a separate backend) produced 1--2 torn-state violations per 25 crash trials at \texttt{crash\_atomicity\_broken.sh}; the honest in-transaction path produced 0 of 25.

%%% ---------------------------------------------------------------
\section{Implementation}
\label{sec:implementation}
%%% ---------------------------------------------------------------

KNDB is a \texttt{shared\_preload\_libraries} extension of 2274 lines of C in \texttt{src/} plus 310 lines of headers in \texttt{include/} (2584 total per \texttt{wc -l}, of which 183 are the F21 C-level test probe \texttt{src/epistemic\_probe.c}, leaving 2401 lines in the production surface). It targets PostgreSQL 18 REL\_18\_STABLE via PGXS. Building the extension against a stock 18.4 install requires no PostgreSQL patch.

\subsection{Handler and delegation}
The AM handler copies the heap AM's \texttt{TableAmRoutine} at first call and overrides seven entries:
\begin{itemize}
\item \texttt{tuple\_insert} (single-row insert path from \texttt{ExecInsert});
\item \texttt{multi\_insert} (batch path from \texttt{COPY FROM CIM\_MULTI}, added at F18);
\item \texttt{tuple\_update} (single-row update path from \texttt{ExecUpdate}, added at F20; rejects any UPDATE that alters \texttt{ep\_kind}, \texttt{ep\_specificity}, or \texttt{ep\_confidence});
\item \texttt{tuple\_delete} (single-row delete path from \texttt{ExecDelete}, added at F20; rejects DELETE outright);
\item \texttt{tuple\_insert\_speculative} (first phase of \texttt{INSERT \dots ON CONFLICT} from \texttt{ExecInsert}, added at F21; runs R1--R5, the advisory lock, and precedence; delegates the speculative write to heap; stashes any pending eviction until \texttt{tuple\_complete\_speculative});
\item \texttt{tuple\_complete\_speculative} (second phase, added at F21; on \texttt{succeeded=true} drains the pending eviction, on \texttt{succeeded=false} discards it so no audit row is written for a killed speculative winner);
\item \texttt{relation\_toast\_am} (delegates TOAST to the same AM so long values stay under the epistemic table's namespace).
\end{itemize}
The remaining $\sim$35 callbacks (\texttt{scan\_begin}, \texttt{scan\_getnextslot}, \texttt{tuple\_lock}, \texttt{index\_fetch\_*}, \texttt{relation\_set\_new\_filelocator}, \texttt{relation\_copy\_data}, \texttt{relation\_vacuum}, and so on) are heap's, unmodified. Reads, index builds, VACUUM, and analyse all inherit heap's behaviour byte-for-byte. Section~\ref{sec:tableam-audit} converts this ``$\sim$35'' claim into an exhaustive audit (Table~\ref{tab:tableam-audit}) with a citable reason per row.

\subsection{The \texttt{tuple\_insert} wrapper}
On each call, \texttt{epistemic\_tuple\_insert\_impl} runs the following sequence:
\begin{enumerate}
\item R1--R5 rule check on the candidate row.
\item Per-slot advisory transaction lock (Section~\ref{sec:design}).
\item \texttt{find\_live\_overlap} scan against \texttt{GetLatestSnapshot()} for a live row in the same (entity\_id, attribute) slot.
\item \texttt{epistemic\_precedence\_cmp} between candidate and incumbent, including the xmin tiebreak.
\item If NEW\_LOSES, raise \texttt{epistemic precedence: NEW\_LOSES (reason=\dots)} and the transaction rolls back.
\item If NEW\_WINS, call \texttt{heapam.tuple\_insert} for the winner, \texttt{epistemic\_close\_sys\_time} for the incumbent, and \texttt{epistemic\_audit\_evicted} for the audit row.
\item Emit one annotation record on custom WAL rmgr 128 (\texttt{epistemic\_wal\_log\_insert\_marker}).
\end{enumerate}
Step 7 is not load-bearing for durability. Section~\ref{sec:design-wal} explains.

\subsection{The \texttt{multi\_insert} wrapper (COPY FROM)}
PostgreSQL 18's \texttt{CopyFrom} at \texttt{copyfrom.c:995-1006} selects \texttt{insertMethod=CIM\_MULTI} when the target has no BEFORE or INSTEAD OF INSERT trigger, then buffers up to 1000-tuple batches (\texttt{MAX\_BUFFERED\_TUPLES}) and flushes them via \texttt{table\_multi\_insert} at \texttt{copyfrom.c:554-559}. Before F18 we copied heap's \texttt{TableAmRoutine} and inherited \texttt{heap\_multi\_insert} unchanged, which meant that every COPY-batched row skipped R1--R5, the advisory lock, precedence, eviction, and the audit. The bad row landed silently; the transcript reads \texttt{COPY 1} or \texttt{COPY 5} with no error.

F18 overrides \texttt{multi\_insert} with a for-loop that invokes \texttt{epistemic\_tuple\_insert\_impl(rel, slots[i], cid, options, bistate)} for each slot in the batch. This is byte-for-byte identical enforcement to a single-row \texttt{INSERT} of the same rows. Correctness is preserved over throughput: \texttt{heap\_multi\_insert} would toast in bulk, pack tuples onto pages with amortised allocation, and emit one \texttt{XLOG\_HEAP2\_MULTI\_INSERT} WAL record per page; the per-slot fan-out pays one \texttt{heap\_insert} per row and one WAL record per row. COPY into an epistemic table therefore runs at roughly the throughput of \texttt{INSERT \dots SELECT}, not of a plain-heap COPY.

We chose full override over loud rejection because \texttt{pg\_dump | pg\_restore} regenerates \texttt{\textbackslash copy} statements and refusing them would break restore for any epistemic table. The alternative to per-slot advisory locks in the batch path would have been to drop the lock; that silently re-opens the RC integrity leak that F6 closed.

\subsection{The \texttt{tuple\_update} and \texttt{tuple\_delete} wrappers (F20)}
\label{sec:impl-f20}
Before F20 the AM inherited \texttt{heapam\_tuple\_update} and \texttt{heapam\_tuple\_delete} verbatim (\texttt{heapam\_handler.c:1384-1385} REL\_18\_STABLE bindings; \texttt{heap\_update} at \texttt{heapam.c:3241}, \texttt{heap\_delete} at \texttt{heapam.c:2772}), so an adversary with \texttt{INSERT} plus \texttt{UPDATE} could overwrite the epistemic prefix on a live row without any rule check, without the advisory lock, without the precedence lattice, and without an eviction audit row --- and an adversary with \texttt{DELETE} could remove a \texttt{MEASURED} incumbent so that a subsequent \texttt{INSERT} landed unopposed. The attack transcripts live in \texttt{scripts/update\_forgery.sh} and \texttt{scripts/delete\_forgery.sh}.

The naive fix --- re-run the \texttt{tuple\_insert} enforcement path on the candidate slot with the incumbent-as-self supplied to the precedence lattice --- fails: a NEW \texttt{MEASURED/1.0} row legitimately beats an OLD \texttt{INFERRED/0.4} row by kind rank, so the attack succeeds through the lattice. The only defensible cut is at the prefix. \texttt{epistemic\_tuple\_update} fetches the incumbent tuple at \texttt{otid} via \texttt{heap\_fetch(rel, GetLatestSnapshot(), \dots)}, deforms it, compares \texttt{(ep\_kind, ep\_specificity, ep\_confidence)} against the candidate slot, and \texttt{ereport(ERRCODE\_CHECK\_VIOLATION)}s if any of the three differ. Legitimate content updates (value, valid\_time, sources) reach \texttt{heapam\_tuple\_update} unchanged, so pt-osm-style column edits still work.

\texttt{epistemic\_tuple\_delete} refuses \texttt{DELETE} outright with \texttt{ERRCODE\_FEATURE\_NOT\_SUPPORTED}. The stricter option --- allow \texttt{DELETE} only for rows whose \texttt{sys\_time} upper bound is closed, and always write an audit row --- would layer atop this and is left to a future eviction API. The rationale for cutting hardest at delete: any \texttt{DELETE} of a live row corresponds to an eviction event that the precedence lattice should have arbitrated, and there is no adversary-friendly workflow that requires the caller to bypass that arbitration.

\subsection{The speculative-insertion wrappers (F21)}
\label{sec:impl-f21}
\texttt{INSERT \dots ON CONFLICT} routes through a two-phase protocol at \texttt{nodeModifyTable.c:1189-1216} REL\_18\_STABLE: \texttt{table\_tuple\_insert\_speculative} writes the tuple with a speculative token, \texttt{ExecInsertIndexTuples} probes the arbiter index for a conflict, and \texttt{table\_tuple\_complete\_speculative} either confirms or kills the tuple. Before F21 we inherited \texttt{heapam\_tuple\_insert\_speculative} and \texttt{heapam\_tuple\_complete\_speculative} verbatim, so an \texttt{ON CONFLICT} that reached the speculative path would skip every enforcement step. As Section~2.3 records, that path is not currently \emph{reachable} from SQL because unique/exclusion index creation fails on epistemic tables at \texttt{heap\_getnext}'s \texttt{rd\_tableam} identity check (\texttt{heapam.c:1352}), so \texttt{ON CONFLICT (col)} has no arbiter to bind. The overrides are defensive; the engine-in-storage claim must not depend on the accidental index-support gap persisting, and the C-level probe at \texttt{sql/am\_speculative.sql} proves they are load-bearing.

\texttt{epistemic\_tuple\_insert\_speculative} runs the same first four steps as the plain insert path: R1--R5, the F6 advisory xact lock, \texttt{find\_live\_overlap}, precedence with the F8 xmin tiebreak. The speculative write itself is delegated to \texttt{heapam->tuple\_insert\_speculative} so heap stamps the tuple with \texttt{HEAP\_INSERT\_SPECULATIVE} and the caller-supplied \texttt{specToken}. Eviction bookkeeping (audit row + \texttt{sys\_time} close + rmgr-128 marker) is deferred: if the arbiter check at \texttt{nodeModifyTable.c:1199} concludes a conflict, \texttt{table\_tuple\_complete\_speculative(succeeded=false)} calls \texttt{heap\_abort\_speculative} (\texttt{heapam.c:6186}) and the winner tuple vanishes. Writing the audit row before the confirm phase would leave the store with an evicted incumbent and no winner --- a durability violation strictly worse than the bypass.

\texttt{epistemic\_tuple\_complete\_speculative} therefore drains a single-entry backend-local pending-eviction slot keyed by \texttt{specToken}. On \texttt{succeeded=true} the slot's contents are folded into \texttt{epistemic\_audit\_evicted} + \texttt{epistemic\_close\_sys\_time} + \texttt{epistemic\_wal\_log\_insert\_marker}, mirroring the plain-insert bookkeeping. On \texttt{succeeded=false} the slot is discarded. The F6 advisory lock is xact-scope and stays held across both callbacks; it is released at outer-transaction commit/abort, not at speculative-complete, so a peer that raced us waits until we finish. The single-entry assumption is safe because \texttt{ExecInsert} at \texttt{nodeModifyTable.c:1189-1216} holds \texttt{SpeculativeInsertionLockAcquire} across both calls, so a backend performs at most one speculative insertion at a time.

\subsection{WAL: annotation, not durability}
\label{sec:design-wal}
The extension registers custom resource manager id 128 (\texttt{RM\_EXPERIMENTAL\_ID}) and emits one \texttt{XLOG\_EPISTEMIC\_INSERT} record per row after the heap insert commits. This record is intentionally minimal: \texttt{tuple\_len=0}, no buffer reference. Heap's \texttt{XLOG\_HEAP\_INSERT} at \texttt{heapam.c:2222-2226} already carries every column of the tuple, including \texttt{ep\_kind}, \texttt{ep\_specificity}, \texttt{ep\_confidence}, and \texttt{sys\_time}, and \texttt{heap\_xlog\_insert} at \texttt{heapam\_xlog.c:482-503} reconstructs the row at redo. We verified empirically (\texttt{scripts/recovery.sh}, 110 rows round-trip; \texttt{wal\_consistency\_checking=all}) that recovery succeeds byte-for-byte with the rmgr-128 emitter disabled. The custom rmgr is retained as a named channel a future logical decoding consumer can hook into; it is not load-bearing for durability of the row.

\subsection{Bulk-load ceiling}
The advisory lock is per-row and lives in PostgreSQL's per-transaction fast-path lock table (\texttt{lock.c:56-57}, \texttt{NLOCKENTS = max\_locks\_per\_xact} $\times$ \texttt{(MaxBackends + max\_prepared\_xacts)}). At PostgreSQL's default \texttt{max\_locks\_per\_transaction=64}, one transaction inserting $\sim$15\,000 distinct-slot rows exhausts the shared lock table and raises \texttt{ERROR 53200: out of shared memory} with the hint to raise \texttt{max\_locks\_per\_transaction}. \texttt{scripts/lock\_exhaustion\_linearity.sh} sweeps the GUC and shows the ceiling scales linearly (raising to 1024 pushes it to $\sim$250\,000). The ceiling applies to COPY too, in the same shape: \texttt{scripts/bypass.sh} sweeps $N \in \{100, 1000, 10000, 20000\}$ via real \texttt{\textbackslash copy} at the default GUC and the first three succeed while $N=20000$ raises \texttt{ERROR 53200} at \texttt{LockAcquireExtended, lock.c:1080}. This is an accepted design tradeoff: dropping the lock in the batch path would re-open the F6 integrity leak; raising the GUC (a one-line change in \texttt{postgresql.conf}) removes the ceiling in a way operators can already control.

\subsection{Threat surface, restated}
Every callback listed above runs from inside \texttt{heapam}. No user-space GUC or \texttt{ALTER TABLE} reaches it. The one exception is \texttt{ALTER TABLE \dots SET ACCESS METHOD heap}, which is a schema-change threat available to the table owner only; it rewrites the relation onto plain heap, at which point the callback is out of the write path entirely, because the callback is no longer bound to the relation. This is out of scope and disclosed in Section~\ref{sec:limitations}.

\subsection{Exhaustive TAM audit}
\label{sec:tableam-audit}
Enumerating write paths anecdotally has been wrong three times in this work: F9 found the \texttt{COPY} \texttt{CIM\_MULTI} bypass, F20 found the \texttt{UPDATE}/\texttt{DELETE} bypass, and F21 found the speculative-insertion bypass. Each was a callback we had classified as ``heap's, unmodified, cannot mutate the epistemic invariant'' and each turned out to route a real write. Table~\ref{tab:tableam-audit} converts anecdotal enumeration into a completeness argument: every one of the 42 callbacks in the PG 18 REL\_18\_STABLE \texttt{TableAmRoutine} interface (\texttt{access/tableam.h}) is classified with a citable file:line reference and a stated reason why it is either overridden or safe to inherit. The classification method is: (i) read the callback signature and docstring; (ii) find heap's binding in \texttt{heapam\_handler.c} and read the body; (iii) grep for every caller in \texttt{src/backend} and confirm the caller either never mutates epistemic-visible state or routes back through a callback we do override. Seven callbacks are \emph{overridden} (each with its own source-rebuild disable-and-test in \texttt{scripts/} and each with the dylib hash flip recorded in \texttt{DECISIONS.md}); the other 35 are inherited. Two rows are flagged \emph{delegated-storage} with the honest disclosure that \texttt{TRUNCATE} and \texttt{CLUSTER}/\texttt{VACUUM FULL} do bypass epistemic policy: both are DDL-privileged, not write-path, so they belong alongside \texttt{SET ACCESS METHOD heap} in the schema-change threat category disclosed in Section~\ref{sec:limitations}.

% tableam_audit.tex --- exhaustive PG 18 REL_18_STABLE TableAmRoutine
% audit. Every callback in access/tableam.h is classified as one of:
%
%   OVERRIDDEN  (7 after F21): our AM binds a wrapper.
%   READ-ONLY   : callback is a scan/fetch/estimate that cannot mutate
%                 epistemic state.
%   DELEGATED-STORAGE: heap's semantics accepted; the AM does not
%                 introduce policy at this level, but any subsequent
%                 mutating path routes through our overridden write
%                 callbacks (e.g., rewrite path repopulates via
%                 tuple_insert).
%   DELEGATED-DECIDE: heap decides on our behalf and we accept the
%                 answer (e.g., TOAST configuration).
%
% Line numbers reference PG 18 REL_18_STABLE
% src/include/access/tableam.h. Every "why" is checked against the
% named citation in the source tree at
% https://git.postgresql.org/gitweb/?p=postgresql.git;a=blob;f=src/include/access/tableam.h;hb=REL_18_STABLE .
%
% See DECISIONS.md (F21) for the classification methodology and the
% disable-and-test proofs behind the seven OVERRIDDEN rows.

\begin{table*}[t]
\centering
\caption{Exhaustive audit of the PG 18 \texttt{TableAmRoutine} interface
  (\texttt{access/tableam.h}, REL\_18\_STABLE, 42 callbacks). The
  seven \textsc{overridden} rows are load-bearing: each has a
  source-rebuild disable-and-test in \texttt{scripts/} that proves the
  bypass returns when the wiring is commented out (see
  \texttt{DECISIONS.md} F9, F18, F20, F21). The other 35 rows are
  inherited from heap; the ``why safe'' column cites the specific
  reason we do not need to intervene.}
\label{tab:tableam-audit}
\scriptsize
\setlength{\tabcolsep}{4pt}
\begin{tabular}{p{3.6cm}p{1.5cm}p{1.5cm}p{9.0cm}}
\toprule
callback & class & tableam.h & why safe / role \\
\midrule
\multicolumn{4}{l}{\emph{Slot callbacks}} \\
\midrule
\texttt{slot\_callbacks}             & read-only          & 302 & Returns the \texttt{TupleTableSlotOps} for tuples of this AM. Pure factory; no state mutation. \\
\midrule
\multicolumn{4}{l}{\emph{Table scan callbacks}} \\
\midrule
\texttt{scan\_begin}                 & read-only          & 326--330 & Starts a read-only scan with a snapshot. No mutation path. \\
\texttt{scan\_end}                   & read-only          & 336 & Releases scan resources. No mutation. \\
\texttt{scan\_rescan}                & read-only          & 342--344 & Restarts scan with new params. No mutation. \\
\texttt{scan\_getnextslot}           & read-only          & 349--351 & Returns next tuple into slot. No mutation. \\
\texttt{scan\_set\_tidrange}         & read-only          & 370--372 & Sets bounds for a TID-range scan. No mutation. \\
\texttt{scan\_getnextslot\_tidrange} & read-only          & 378--380 & TID-range variant of \texttt{scan\_getnextslot}. No mutation. \\
\midrule
\multicolumn{4}{l}{\emph{Parallel scan}} \\
\midrule
\texttt{parallelscan\_estimate}      & read-only          & 391 & Sizes shared-mem DSM. No mutation. \\
\texttt{parallelscan\_initialize}    & read-only          & 398--399 & Initialises a \texttt{ParallelTableScanDesc}. No user-visible state. \\
\texttt{parallelscan\_reinitialize}  & read-only          & 405--406 & Re-initialises the DSM for a new scan. No mutation. \\
\midrule
\multicolumn{4}{l}{\emph{Index scan}} \\
\midrule
\texttt{index\_fetch\_begin}         & read-only          & 422 & Opens an index-scan fetch state. No mutation. \\
\texttt{index\_fetch\_reset}         & read-only          & 428 & Resets between index scans. No mutation. \\
\texttt{index\_fetch\_end}           & read-only          & 433 & Releases index-scan state. No mutation. \\
\texttt{index\_fetch\_tuple}         & read-only          & 455--459 & Fetches by TID after visibility test. No mutation. \\
\midrule
\multicolumn{4}{l}{\emph{Non-modifying tuple ops}} \\
\midrule
\texttt{tuple\_fetch\_row\_version}  & read-only          & 472--475 & Fetches a tuple version. No mutation. \\
\texttt{tuple\_tid\_valid}           & read-only          & 480--481 & Predicate on tid. No mutation. \\
\texttt{tuple\_get\_latest\_tid}     & read-only          & 487--488 & Walks the update chain to the newest visible tid. Reads only. \\
\texttt{tuple\_satisfies\_snapshot}  & read-only          & 494--496 & Visibility test. No mutation. \\
\texttt{index\_delete\_tuples}       & delegated-storage  & 499--500 & Bottom-up index-delete driver. In heap's implementation, no user-visible tuples are removed; dead line pointers are reclaimed. Our epistemic invariants live in \emph{live} rows (\texttt{sys\_time @> now()}); dead-line-pointer reclamation cannot affect them. \\
\midrule
\multicolumn{4}{l}{\emph{Mutating tuple ops}} \\
\midrule
\texttt{tuple\_insert}               & \textbf{OVERRIDDEN} & 508--511 & F1--F16: R1--R5, F6 advisory lock, precedence + F8 xmin tiebreak, eviction audit, sys\_time close, rmgr-128 marker. \\
\texttt{tuple\_insert\_speculative}  & \textbf{OVERRIDDEN} & 513--519 & \textbf{F21}: R1--R5, advisory lock, precedence; heap does the speculative write; deferred eviction stashed for confirm. \\
\texttt{tuple\_complete\_speculative}& \textbf{OVERRIDDEN} & 521--525 & \textbf{F21}: on succeeded=true, drain the pending eviction (audit + sys\_time close + rmgr-128); on succeeded=false, discard. \\
\texttt{multi\_insert}               & \textbf{OVERRIDDEN} & 527--529 & F18: fan out to \texttt{epistemic\_tuple\_insert\_impl} per slot; closes the COPY CIM\_MULTI bypass F9 identified. \\
\texttt{tuple\_delete}               & \textbf{OVERRIDDEN} & 531--539 & F20: refuse \texttt{DELETE} outright with \texttt{ERRCODE\_FEATURE\_NOT\_SUPPORTED}. \\
\texttt{tuple\_update}               & \textbf{OVERRIDDEN} & 541--551 & F20: refuse any \texttt{UPDATE} whose new \texttt{(ep\_kind, ep\_specificity, ep\_confidence)} differs from the incumbent's. \\
\texttt{tuple\_lock}                 & delegated-storage  & 553--562 & Row-level lock acquisition (\texttt{SELECT FOR UPDATE}, replica-apply conflict resolution). Sets \texttt{xmax} to lock the tuple but does not mutate user columns; any subsequent modification routes back through \texttt{tuple\_update} or \texttt{tuple\_delete}, both overridden. \\
\texttt{finish\_bulk\_insert}        & delegated-storage  & 576 & Optional per-bulk-insert flush. In-tree AMs no longer use it; heap's binding is \texttt{NULL} (\texttt{heapam\_handler.c:2661}). Any prior mutation went through \texttt{tuple\_insert} or \texttt{multi\_insert}. \\
\midrule
\multicolumn{4}{l}{\emph{DDL}} \\
\midrule
\texttt{relation\_set\_new\_filelocator} & delegated-storage & 600--604 & Allocates a new physical file for TRUNCATE/CLUSTER/REINDEX-style rewrites (\texttt{heapam\_handler.c:583--622}: \texttt{RelationCreateStorage} + optional init-fork). Storage is empty on return; any subsequent rewrite repopulates via \texttt{tuple\_insert} or \texttt{multi\_insert}. \\
\texttt{relation\_nontransactional\_truncate} & delegated-storage & 614 & Truncates the file to zero size (heap: \texttt{RelationTruncate}, no WAL). Removes all rows including epistemic incumbents. This IS a policy-relevant bypass surface for the eviction audit --- but it is \texttt{TRUNCATE} DDL, guarded by \texttt{TRUNCATE} privilege, and out of scope for the write-path threat (Section~2.3). Flagged as a schema-change bypass alongside \texttt{ALTER TABLE \dots SET ACCESS METHOD heap} (Section~9). \\
\texttt{relation\_copy\_data}        & delegated-storage & 622--623 & Copies the physical file to a new tablespace (\texttt{ALTER TABLE \dots SET TABLESPACE}). Rows are copied byte-for-byte; no epistemic-level policy applies. \\
\texttt{relation\_copy\_for\_cluster}& delegated-storage & 626--635 & Copies rows during \texttt{CLUSTER} / \texttt{VACUUM FULL}. Heap re-inserts via \texttt{raw\_heap\_insert}, not \texttt{table\_tuple\_insert}, so the epistemic checks do NOT re-fire during the rewrite. Accepted because CLUSTER/VACUUM FULL is a DDL-privileged operation that preserves the current visible rows; it cannot introduce a forgery, only reorder existing rows. Flagged as an in-scope limitation on par with the schema-change threat. \\
\texttt{relation\_vacuum}            & delegated-storage & 652--654 & Standard VACUUM (dead-tuple reclamation, freeze). Cannot resurrect a dead row or invent a live one; the ``live row per slot'' invariant is preserved. \\
\texttt{scan\_analyze\_next\_block}  & read-only          & 673--674 & ANALYZE block sampling. Read-only. \\
\texttt{scan\_analyze\_next\_tuple}  & read-only          & 684--688 & ANALYZE tuple sampling. Read-only. \\
\texttt{index\_build\_range\_scan}   & read-only          & 691--701 & Feeds tuples to \texttt{IndexBuildCallback} during index build. Read-only from the AM's perspective; the index AM writes to its own relation. \\
\texttt{index\_validate\_scan}       & read-only          & 704--708 & Concurrent-index-build validation phase. Read-only against the table. \\
\midrule
\multicolumn{4}{l}{\emph{Misc}} \\
\midrule
\texttt{relation\_size}              & read-only          & 724 & Returns bytes for a fork. Pure metadata. \\
\texttt{relation\_needs\_toast\_table} & delegated-decide & 734 & Heap's decision, whether a TOAST table is needed. We accept it. \\
\texttt{relation\_toast\_am}         & \textbf{OVERRIDDEN} & 741 & Return \texttt{HEAP\_TABLE\_AM\_OID} so the TOAST table is plain heap; otherwise PG's index build for the TOAST table trips \texttt{heap\_getnext}'s \texttt{rd\_tableam == GetHeapamTableAmRoutine()} identity check (\texttt{heapam.c:1352}). \\
\texttt{relation\_fetch\_toast\_slice} & read-only        & 748--752 & Fetches a slice of a TOAST'd value. Read-only. \\
\midrule
\multicolumn{4}{l}{\emph{Planner}} \\
\midrule
\texttt{relation\_estimate\_size}    & read-only          & 770--772 & Row-count / page-count estimate for the planner. Pure. \\
\midrule
\multicolumn{4}{l}{\emph{Executor}} \\
\midrule
\texttt{scan\_bitmap\_next\_tuple}   & read-only          & 792--796 & Fetch next visible tuple from a bitmap scan. Read-only. \\
\texttt{scan\_sample\_next\_block}   & read-only          & 823--824 & Sample-scan block selection. Read-only. \\
\texttt{scan\_sample\_next\_tuple}   & read-only          & 839--841 & Sample-scan tuple selection. Read-only. \\
\bottomrule
\end{tabular}
\end{table*}


%%% ---------------------------------------------------------------
\section{Formalism}
\label{sec:formalism}
%%% ---------------------------------------------------------------

The Sybil-collapse threshold of an agreement-based truth-discovery algorithm is predictable from the dataset's per-slot honest-support distribution alone. This section makes that statement precise and validates it empirically on the two datasets in the evaluation. The complete derivation is in the companion \texttt{bench/docs/sybil\_formalism.md}; this section reports what the paper needs.

\subsection{Model and threat}
Let $S = S_h \cup S_s$ be the disjoint sets of honest and Sybil sources, $I$ the items, $v^{*}(i)$ the true value at item $i$. A claim is a triple $(i, s, v)$. Define
\[
h_{\text{top}}(i) = \max_{v \in V^{*}_{i}} \left| \{ s \in S_h : (i, s, v) \text{ claimed} \} \right|,
\]
where $V^{*}_{i}$ is the set of value strings the correctness scorer accepts as matching gold. For binary tasks (Zheng d\_sentiment) $V^{*}_{i} = \{v^{*}(i)\}$ and $h_{\text{top}}(i) = h(i)$. For string-valued tasks (Book-Author) gold-matching claims are split across surface forms (\texttt{"O'Leary, Timothy J."} versus \texttt{"Timothy J O'Leary"}), so $h_{\text{top}}(i) \leq h(i)$.

Each TD algorithm assigns each source a weight $w(s) \geq 0$ and picks
$
v_{\hat{}}(i) = \arg\max_v \sum_{s : c(s,i)=v} w(s)
$
with algorithm-specific tie-break. A Sybil coalition of size $k$ per slot asserts one falsified value $f(i) \neq v^{*}(i)$.

\subsection{Per-algorithm monotonicity}
The four algorithms in the evaluation are TruthFinder~\cite{yin2007truthfinder}, CRH~\cite{li2014crh}, CATD~\cite{li2015catd}, and ACCU~\cite{dong2009accu}. For each, the summed weight of $k$ Sybils asserting one value grows monotonically in $k$: TruthFinder's fixed-point iteration bootstraps supporter trust from mutual agreement (KDD 2007 eqs.\ 3, 7, 8); CRH's log-ratio yields a defensive $O(1/\ln|I|)$ margin against a single Sybil but not against a coalition; CATD's $\chi^{2}_{\alpha/2, n}$ upper confidence bound tightens as $n$ grows but the summed weight is linear in $k$; ACCU's Bayesian MAP over per-source accuracy has a per-supporter log-odds coefficient that stays positive for any $A > 1/(1+n_{\text{false}})$. Details in the companion document.

\subsection{Threshold theorem}
Under standard hyperparameters (TruthFinder: $\gamma=0.3$, $\rho=0.5$, initial trust 0.9; CRH: default; CATD: $\alpha=0.05$; ACCU: initial $A=0.8$), for each of the four algorithms there exists a finite threshold
\[
k^{*}(i) \;\lesssim\; \rho_{\text{alg}} \cdot h_{\text{top}}(i)
\]
above which the Sybil coalition flips $v_{\hat{}}(i)$ from $v^{*}(i)$. The per-algorithm constants $\rho$ are bounded above by $1 + o(1)$: $\rho_{\text{TF}} \approx 1$ (sometimes below 1 in the sub-saturation regime because the iteration is self-amplifying); $\rho_{\text{CRH}} \approx 1 + 1/\ln|I|$; $\rho_{\text{CATD}} \approx 1$ at $\alpha=0.05$; $\rho_{\text{ACCU}} \approx 1$ on binary tasks. The consequence is $k^{*} = \Theta(h_{\text{top}})$: the Sybil-collapse cell is linear in per-slot honest support with a per-algorithm slope close to one.

\subsection{KNDB invariance}
Under the F1--F8 KNDB lattice with tier mapping computed from independent metadata, $k^{*}$ is unbounded. Proof: the lattice orders \texttt{kind} strictly above every other axis, and \texttt{kind $\in \{$MEASURED, INFERRED, DERIVED$\}$} is assigned at write time by a rule that consults only metadata living in the dataset before conflict resolution. On Book-Author, Tier A (MEASURED) requires membership in the top $K/2$ by \texttt{n\_listings} \emph{and} \texttt{canon\_rate} $\geq 0.5$; an adversarial identity appearing for the first time in the trace has \texttt{n\_listings} $\leq k \ll K/2$ and no \texttt{canon\_rate} history. On Zheng, Tier A requires top-third \texttt{quali\_acc} on the disjoint qualification test (question IDs 2000--2019); an adversarial identity has no qualification responses. In both cases the adversary is at most INFERRED. MEASURED strictly outranks INFERRED, so a single Tier-A MEASURED honest write beats any coalition of INFERRED writes regardless of coalition size or confidence values. Therefore $k^{*}(i) = \infty$.

The empirical counterpart of this theorem is the F14/F15b KIND\_OFF disable-and-test (Section~\ref{sec:evaluation}): when the kind-primary ordering is patched out and only confidence remains, KNDB collapses to \texttt{pg\_conf}'s numeric behaviour on both datasets.

\subsection{Empirical validation}
We measured $h_{\text{top}}$ directly on both datasets. Book-Author K=50: median $h_{\text{top}} = 4$, mean 7.3, p90 18. Predicted collapse bracket $k^{*} \in [3, 5]$; observed sharp cliff at N=5\ldots10 for TruthFinder, CATD, and ACCU (TruthFinder $0.530 \to 0.120$ at N=5, $\to 0.010$ at N=10; CATD $0.550 \to 0.420 \to 0.100$; ACCU $0.530 \to 0.290 \to 0.060$); CRH lags one step ($0.580$ at N=5, $0.230$ at N=10). The predicted bracket is within $\pm 20\%$ of the observed cliff.

Zheng d\_sentiment K=45: median $h_{\text{top}} = 14$, uniform 20 votes per slot. Predicted collapse bracket $k^{*} \in [14, 20]$; observed collapse at N=20 with CRH, CATD, and ACCU all going to zero. Within $\pm 30\%$ of the observed threshold on the $k$-axis; the observed cell is the density-saturation cell exactly.

%%% ---------------------------------------------------------------
\section{Evaluation}
\label{sec:evaluation}
%%% ---------------------------------------------------------------

\subsection{Experimental setup}
All measurements are on PostgreSQL 18.4 at \texttt{/tmp/kndb\_pg18\_test:55480}, one Apple M5 Pro (18 cores, 48 GB RAM, macOS 26.4.1), unix-socket connections. The extension is loaded via \texttt{shared\_preload\_libraries = 'epistemic'} and defaults otherwise, except where a specific cell required raising \texttt{max\_locks\_per\_transaction} for a bulk-load probe. Every bench cell records the extension dylib SHA-256 in its metadata block. The reference honest build in this paper is dylib SHA-256 \texttt{eb15d442dd1eace5\-88c1c4ee4fab3e18\-3addc70cc3a28773\-f8d5acfc0ff58af0}, the F18 post-COPY-close artefact. Cells taken before F18 are on \texttt{807b2e87\ldots}, disclosed inline where relevant. Every source-rebuild disable-and-test recorded both the honest and patched dylib SHA-256 and confirmed the source tree returned byte-identical to HEAD after the run.

\subsection{Confidence-forgery workload (Book-Author)}
Dataset: Dong et al.\ VLDB 2009 Book-Author fusion~\cite{dong2009accu}, 895 bookstores, 1265 books, 33\,971 assertions, gold for 100 ISBNs. Tier map (independence audit in \texttt{bench/datasets/bookauthor/README.md}): Tier A (MEASURED, conf uniform $[0.5, 0.9]$) if top-$K/2$ by \texttt{n\_listings} \emph{and} \texttt{canon\_rate} $\geq 0.5$; Tier B (INFERRED, conf $[0.4, 0.7]$); Tier C (DERIVED, conf $[0.2, 0.5]$). Adversarial injection ($N$ per gold ISBN) is INFERRED with conf $[0.95, 1.0]$ carrying a scrambled real author name from a different gold ISBN. K=50 held fixed; N swept over $\{1, 3, 5, 10\}$; seed 20260714 fixes RNG.

Table~\ref{tab:bookauthor} reports precision on gold ISBNs at $c=1$. KNDB stays at 0.630 across every N (kind rank picks Tier-A MEASURED over adversarial INFERRED); \texttt{pg\_conf} collapses to 0.000 at every N (ranked by confidence alone, the adversarial $[0.95, 1.0]$ beats Tier-A $[0.5, 0.9]$); \texttt{pg\_lww} degrades $0.210 \to 0.040$; \texttt{pg\_mv} $0.460 \to 0.170$; \texttt{pg\_heap} is an integrity failure (max 124 live rows per slot, mean 38.6) and its numeric precision (0.68--0.71) is a coin flip over which duplicate the scan returned first. \texttt{pg\_trigger} tracks KNDB to within one point (same lattice, plpgsql implementation, marginally higher abort rate under contention).

\begin{table}[t]
\centering
\caption{Book-Author confidence forgery, precision on 100 gold ISBNs, $c=1$. From \texttt{bench/results/summary/stage3\_adversarial.md}.}
\label{tab:bookauthor}
\small
\begin{tabular}{lcccc}
\toprule
system & N=1 & N=3 & N=5 & N=10 \\
\midrule
KNDB epistemic & 0.630 & 0.630 & 0.630 & 0.630 \\
pg\_trigger    & 0.620 & 0.620 & 0.620 & 0.620 \\
pg\_mv         & 0.460 & 0.400 & 0.320 & 0.170 \\
pg\_lww        & 0.210 & 0.130 & 0.090 & 0.040 \\
pg\_conf       & 0.000 & 0.000 & 0.000 & 0.000 \\
pg\_heap       & \multicolumn{4}{c}{INTEGRITY FAIL (max=124, mean=38.6)} \\
\bottomrule
\end{tabular}
\end{table}

\subsection{Confidence-forgery workload (Zheng)}
Dataset: Zheng et al.\ VLDB 2017 d\_sentiment crowdsourcing task~\cite{zheng2017crowdsourcing}, 85 AMT workers, 1000 items, $\sim$20 labels per item, 999 of 1000 items contested. Tier map (independence audit in \texttt{bench/datasets/zheng\_sentiment/README.md}): rank workers by \texttt{quali\_acc(w)} on the disjoint qualification test (item IDs 2000--2019, 1700 responses); Tier A (MEASURED) = top $K/3$; Tier B (INFERRED); Tier C (DERIVED, all workers outside top-K). Adversarial injection ($N$ per contested slot) is INFERRED with conf $[0.95, 1.0]$ and value = binary flip of gold. K=45 held fixed. Table~\ref{tab:zheng} reports precision at $c=1$.

\begin{table}[t]
\centering
\caption{Zheng d\_sentiment confidence forgery, precision, $c=1$. From \texttt{bench/results/summary/stage3\_zheng\_adversarial.md}.}
\label{tab:zheng}
\small
\begin{tabular}{lcccc}
\toprule
system & N=1 & N=3 & N=5 & N=10 \\
\midrule
KNDB epistemic & 0.927 & 0.927 & 0.927 & 0.927 \\
pg\_trigger    & 0.927 & 0.927 & 0.927 & 0.927 \\
pg\_lww        & 0.401 & 0.211 & 0.130 & 0.074 \\
pg\_mv         & 0.375 & 0.193 & 0.127 & 0.070 \\
pg\_conf       & 0.000 & 0.000 & 0.000 & 0.000 \\
pg\_heap       & \multicolumn{4}{c}{INTEGRITY FAIL} \\
\bottomrule
\end{tabular}
\end{table}

The delta over \texttt{pg\_conf} is 92.7 points, flat across N. The delta over KNDB's second-best PostgreSQL baseline (\texttt{pg\_lww}) is 52--85 points depending on N.

\subsection{Source-rebuild disable-and-test}
On both datasets we patched \texttt{epistemic\_precedence\_cmp} (\texttt{src/epistemic\_rules.c:379-423}) to force \texttt{inc\_rank = new\_rank = 1}, short-circuiting the kind branch so the function falls through to specificity and confidence. This is the exact ranking behaviour of \texttt{pg\_conf} on both workloads (specificity is 0 across the board; confidence decides). After each measurement the source was restored byte-identical (\texttt{git diff --stat contrib/epistemic/src/} empty), rebuilt and reinstalled, and the dylib hash returned to the honest value.

Book-Author, N=5, $c=1$: KNDB kind ON precision 0.630, kind OFF precision 0.000 (dylib flip \texttt{807b2e87\ldots} $\to$ f7\ldots\ $\to$ \texttt{807b2e87\ldots}). Delta $-63$ points, exactly the KNDB-versus-\texttt{pg\_conf} margin. Zheng, N=5, $c=1$: KNDB kind ON precision 0.927, kind OFF precision 0.000 (dylib flip \texttt{807b2e87\ldots} $\to$ \texttt{3cc4f4b7\ldots} $\to$ \texttt{807b2e87\ldots}). Delta $-92.7$ points, exactly the KNDB-versus-\texttt{pg\_conf} margin.

Integrity held under the patched build in both cases (no slots with more than one live row) because the F6 advisory lock and F8 xmin tiebreak still operated; only kind-rank decision-making was disabled. This is the correct decomposition: integrity and correctness are separable mechanisms in KNDB, and each has its own disable-and-test.

\subsection{UPDATE and DELETE forgery attacks (F20)}
\label{sec:eval-updateDelete}
Two attacks the earlier bypass-table enumeration missed: an \texttt{UPDATE} that rewrites the epistemic prefix on a committed row, and a \texttt{DELETE} that removes an incumbent so that a subsequent \texttt{INSERT} lands unopposed. Both were reproduced against a fresh cluster before writing the fix: \texttt{scripts/update\_forgery.sh} seeds an \texttt{INFERRED/0.4} row, issues \texttt{UPDATE fact SET ep\_kind='MEASURED', ep\_confidence=1.0, value='forged'}, and reads back \texttt{MEASURED/1.0/forged} without any error, without any eviction audit row. \texttt{scripts/delete\_forgery.sh} seeds a \texttt{MEASURED} incumbent, deletes it, inserts an \texttt{INFERRED/0.99} row, and reads back \texttt{INFERRED/0.99/benign} as the live row for the slot.

F20 adds two callbacks. \texttt{epistemic\_tuple\_update} fetches the incumbent at \texttt{otid} via \texttt{heap\_fetch} with \texttt{GetLatestSnapshot}, deforms it, and refuses any \texttt{UPDATE} whose new \texttt{(ep\_kind, ep\_specificity, ep\_confidence)} triple differs from the incumbent's, raising \texttt{ERRCODE\_CHECK\_VIOLATION}. \texttt{epistemic\_tuple\_delete} refuses \texttt{DELETE} outright with \texttt{ERRCODE\_FEATURE\_NOT\_SUPPORTED}. Content-only \texttt{UPDATE}s that change value, valid\_time, or sources delegate to \texttt{heapam\_tuple\_update} unchanged, so a legitimate content-correction workflow still succeeds. The regression test cell \texttt{sql/am\_update\_delete.sql} exercises both callbacks under \texttt{make installcheck}.

The disable-and-test proof: with the F20 wiring in \texttt{epistemic\_am\_handler} commented out and the extension rebuilt (dylib flip \texttt{c873ddc4\ldots} $\to$ \texttt{1fb0d572\ldots}), both forgery scripts print \texttt{FORGERY SUCCEEDED}; restoring the two wiring lines byte-identical and rebuilding (\texttt{1fb0d572\ldots} $\to$ \texttt{c873ddc4\ldots}) returns both to \texttt{FORGERY REJECTED}. This confirms the wiring is load-bearing and neither callback is dead code, mirroring the F18 discipline for \texttt{multi\_insert}. Under the honest build, the extension's \texttt{make check-e2e} suite now passes 10/10 (adding \texttt{check-e2e-updforge} and \texttt{check-e2e-delforge} to the F19 suite of 8), and \texttt{scripts/bypass.sh} passes 28 assertions covering five bypass mechanisms $\times$ two tables plus the F7 lock-table sweep.

\subsection{Speculative-insertion attack (F21)}
\label{sec:eval-f21}
F21 adds two callbacks. \texttt{epistemic\_tuple\_insert\_speculative} runs R1--R5, the F6 advisory xact lock, and precedence (with the F8 xmin tiebreak), then delegates the actual write to heap's \texttt{heapam\_tuple\_insert\_speculative} so the row is stamped with \texttt{HEAP\_INSERT\_SPECULATIVE} and the caller's \texttt{specToken}. Eviction bookkeeping is deferred to \texttt{epistemic\_tuple\_complete\_speculative}: on \texttt{succeeded=true} the callback drains a backend-local pending-eviction slot (audit row + \texttt{sys\_time} close + rmgr-128 marker); on \texttt{succeeded=false} it discards the slot, because \texttt{heap\_abort\_speculative} has just removed the winner tuple. The two-phase deferral is load-bearing: writing the audit row before the confirm phase would leave the store with an evicted incumbent and no winner if the arbiter check aborted the speculative row --- worse than the bypass we are closing.

The SQL-level attack surface for the speculative path is not currently reachable on epistemic tables. \texttt{CREATE UNIQUE INDEX}, \texttt{ADD PRIMARY KEY}, \texttt{ADD UNIQUE}, and \texttt{ADD EXCLUDE} all fail on epistemic relations at \texttt{heap\_getnext}'s \texttt{rd\_tableam} identity check (\texttt{heapam.c:1352} REL\_18\_STABLE) during the btree/gist ambuild scan, so \texttt{ON CONFLICT (col)} has no arbiter to bind. \texttt{scripts/speculative\_forgery.sh} confirms empirically that four SQL-level attempts (\texttt{DO UPDATE} with a UNIQUE arbiter that never gets built, \texttt{DO NOTHING} with a target column list, bare \texttt{DO NOTHING}, and constraint-creation probes) all either fail at DDL or fall back to a route already covered by the plain-\texttt{tuple\_insert} override. That is a happy accident of the same index-support gap that drives the $\sim$71\% overhead in Section~\ref{sec:evaluation}; the engine-in-storage claim must not depend on it.

The C-level probe \texttt{epistemic.\_probe\_speculative\_insert}, added at F21 alongside the two overrides, invokes \texttt{table\_tuple\_insert\_speculative} and \texttt{table\_tuple\_complete\_speculative} programmatically with a caller-supplied candidate tuple, mirroring \texttt{ExecInsert}'s two-phase call sequence at \texttt{nodeModifyTable.c:1189-1216}. The regression cell \texttt{sql/am\_speculative.sql} exercises the probe against seven distinct cases: valid MEASURED, R3 violation, valid INFERRED, R4 violation, precedence NEW\_LOSES, precedence NEW\_WINS with deferred eviction, and speculative-abort with a would-be eviction that must be discarded. The disable-and-test proof: with the F21 wiring in \texttt{epistemic\_am\_handler} commented out and the extension rebuilt (dylib flip \texttt{959d5e67\ldots} $\to$ \texttt{c44b276d\ldots}), the probe with an R3-violating MEASURED row and the probe with an R4-violating INFERRED row both return \texttt{OK} and the forged row lands in the target relation; restoring the two wiring lines byte-identical and rebuilding (\texttt{c44b276d\ldots} $\to$ \texttt{959d5e67\ldots}) returns both to the expected \texttt{ERRCODE\_CHECK\_VIOLATION}. Under the honest build, \texttt{make installcheck} passes 8/8 (adding \texttt{am\_speculative} to the F20 suite of 7).

\subsection{Truth-discovery baselines}
We reimplemented TruthFinder, CRH, CATD, and ACCU in Python 3 from the original equations (\texttt{bench/scripts\_td/td\_algorithms.py}). No vendored code. Each algorithm reproduced the Zheng VLDB 2017 survey's D\_PosSent baseline to within 0.3 percentage points (CATD 0.957 vs.\ survey 0.960; CRH 0.950 vs.\ survey PM 0.9504). Default hyperparameters per each paper; no tuning either way. TD algorithms consume the same normalised trace KNDB and \texttt{pg\_*} consume; no trace edits, no MEASURED-row filtering.

Book-Author Sybil at N=10, $c=1$ (Table~\ref{tab:td-sybil}): KNDB 0.630 flat while all four TD algorithms collapse (TruthFinder 0.010, CRH 0.230, CATD 0.100, ACCU 0.060). Best TD (CRH) beaten by 40 points; worst (TruthFinder) by 62 points. Independent-value attack (each adversarial identity picks a distinct wrong value, no coordination): TD is flat across N=1\ldots 10, and the KNDB win narrows to 5--10 points, because independent adversaries get no trust bootstrap.

\begin{table}[t]
\centering
\caption{Book-Author Sybil coordination, precision, $c=1$. TD baselines reproduce Zheng VLDB 2017 survey within 0.3pp on their standard honest baseline. From \texttt{bench/results/summary/stage3\_td\_baselines.md}.}
\label{tab:td-sybil}
\small
\begin{tabular}{lcccc}
\toprule
system & N=1 & N=3 & N=5 & N=10 \\
\midrule
KNDB epistemic & 0.630 & 0.630 & 0.630 & 0.630 \\
CRH            & 0.580 & 0.590 & 0.590 & 0.230 \\
CATD           & 0.550 & 0.550 & 0.420 & 0.100 \\
ACCU           & 0.530 & 0.530 & 0.290 & 0.060 \\
TruthFinder    & 0.530 & 0.470 & 0.120 & 0.010 \\
pg\_mv         & 0.490 & 0.390 & 0.330 & 0.260 \\
pg\_conf       & 0.000 & 0.000 & 0.000 & 0.000 \\
\bottomrule
\end{tabular}
\end{table}

\subsection{Zheng Sybil sweep and density saturation}
Table~\ref{tab:zheng-sybil} shows a broader Sybil sweep on Zheng d\_sentiment, including N=20 (density saturation, Sybil count equals total honest votes per slot). KNDB stays flat at 0.927. CRH holds at 0.951 through N=10 and then drops to 0.000 at N=20 (a sharp cliff, not a gradual degradation). CATD holds at 0.948 through N=5, drops to 0.000 at N=10, stays at 0.000 at N=20. ACCU peaks at 1.000 at N=5 (above KNDB), drops to 0.000 at N=10 and N=20. TruthFinder degrades gradually from 0.905 down to 0.482.

\begin{table}[t]
\centering
\caption{Zheng d\_sentiment Sybil sweep, precision, $c=1$. N=20 is the density-saturation cell. From \texttt{bench/results/summary/stage3\_zheng\_sybil.md}.}
\label{tab:zheng-sybil}
\small
\begin{tabular}{lccccc}
\toprule
system & N=1 & N=3 & N=5 & N=10 & N=20 \\
\midrule
KNDB epistemic & 0.927 & 0.927 & 0.927 & 0.927 & 0.927 \\
CRH            & 0.953 & 0.951 & 0.951 & 0.951 & 0.000 \\
CATD           & 0.955 & 0.953 & 0.948 & 0.000 & 0.000 \\
ACCU           & 0.964 & 0.997 & 1.000 & 0.000 & 0.000 \\
TruthFinder    & 0.905 & 0.690 & 0.557 & 0.494 & 0.482 \\
\bottomrule
\end{tabular}
\end{table}

The disable-and-test on the TD side (replacing each algorithm's weight-update loop with plain majority vote) shows the mechanism is bimodal, not uniformly amplifying. At sub-saturation N=1\ldots 5, CRH, CATD, and ACCU each beat plain MV by 4--26 percentage points (their log-ratio, confidence bound, or MAP identifies wrong Sybils correctly). At N=10, CATD and ACCU flip and score below MV by 29 points (the mechanism inverts: agreement now amplifies rather than defends). At N=20 all three collapse to the MV floor. TruthFinder is the only algorithm that amplifies at low N. This bimodality is the ``TD is not uniformly bad'' finding that shapes the Related Work discussion.

\subsection{Integrity column}
Table~\ref{tab:integrity} reports the integrity axis on the Stage-2 YCSB grid. The grid runs three kind mixes, two contention levels (theta), and two concurrencies (c=8, c=32) against seven systems, for 78 cells total: six systems (KNDB, pg\_heap, pg\_conf, pg\_lww, pg\_mv, pg\_trigger) at $3 \times 2 \times 2 = 12$ cells each, and pg\_llm capped at 6 cells because the calibrated 1120 ms per-conflict latency makes higher-concurrency pg\_llm runs impractical. A cell is INTEGRITY FAIL if any (entity, attribute) slot ended the trace with more than one live row. KNDB is the only system that passes on every cell.

\begin{table}[t]
\centering
\caption{Integrity FAIL counts on the Stage-2 grid (12 cells per system; \texttt{pg\_llm} at 6). From \texttt{bench/results/summary/stage2.md} F14 addendum.}
\label{tab:integrity}
\small
\begin{tabular}{lc}
\toprule
system & INTEGRITY FAIL cells \\
\midrule
KNDB epistemic & 0 / 12 \\
pg\_conf       & 12 / 12 \\
pg\_heap       & 12 / 12 \\
pg\_lww        & 10 / 12 \\
pg\_trigger    & 10 / 12 \\
pg\_mv         & 9 / 12 \\
pg\_llm        & 3 / 6 \\
\bottomrule
\end{tabular}
\end{table}

Every trigger-based baseline fails integrity at 32 clients on the hotter theta values, because two concurrent writers each find ``no incumbent'' via \texttt{SELECT FOR UPDATE} on an empty result set, both commit, and two live rows land. The failure is not unique to \texttt{pg\_lww}; \texttt{pg\_conf}, \texttt{pg\_trigger}, \texttt{pg\_mv}, and \texttt{pg\_llm} all exhibit it at high enough contention. KNDB's F6 advisory lock is what closes this window, and \texttt{scripts/rc\_invariant.sh} confirms the closure by adversarial rebuild.

\subsection{Real LLM calibration and non-determinism}
The mock LLM in the Stage-2 baseline is calibrated to a real Anthropic Claude Haiku 4.5 measurement (\texttt{bench/results/stage3\_llm\_calibration.jsonl}, 200 real API calls, zero errors, model \texttt{claude-haiku-4-5}~\cite{anthropic2025haiku}). Correctness 92.5\% (185 of 200); latency mean 1120 ms, p50 940 ms, p95 1934 ms, p99 2312 ms, max 5597 ms. The mock's earlier assumed correctness of 0.65 and latency of 300 ms were both wrong; parameters were updated to match the measurement.

A separate non-determinism probe (\texttt{bench/results/stage3\_llm\_nondeterminism\_raw.jsonl}, 500 real API calls: 50 conflicts $\times$ 10 replays each) measured a flip rate of 0 of 50 conflicts (all unanimous over the 10 replays) with 2 of 50 unanimous-but-wrong. Both wrong conflicts had the same signature: incumbent INFERRED with specificity 0 and confidence 0.5, candidate INFERRED with high specificity ($\gg 0$) and confidence below 0.5. KNDB's lattice picks the candidate (higher specificity within the same kind). The LLM picks the incumbent every time: it treats ``same kind, higher confidence'' as decisive over ``same kind, higher specificity,'' which is the opposite of the lattice's (spec, conf) precedence order. This is a systematic disagreement about lattice ordering, not random noise; the ``just run the LLM three times and vote'' fix does not help.

\subsection{Bypass table}
Table~\ref{tab:bypass} summarises the F2, F18, and F20 bypass-survival results. Each cell corresponds to a scenario in \texttt{scripts/bypass.sh}, \texttt{scripts/update\_forgery.sh}, or \texttt{scripts/delete\_forgery.sh}, which assert against both an epistemic table (\texttt{fact\_native}) and a plain-heap table with a BEFORE-INSERT/UPDATE/DELETE trigger (\texttt{fact\_trigger}). The adversarial disable-and-test rebuilds the extension with the corresponding TAM callback wiring line commented out and confirms the bypass returns; on F20 the dylib SHA-256 flips \texttt{c873ddc4\ldots} $\to$ \texttt{1fb0d572\ldots} $\to$ \texttt{c873ddc4\ldots}. The trigger column marks a scenario \emph{bypassed} whenever a writer with \texttt{INSERT}+\texttt{ALTER TABLE} on the target can turn the trigger off; \emph{bypassed*} marks the two rows where the trigger-based enforcer would survive if the operator opts in to \texttt{ENABLE ALWAYS TRIGGER} but still cannot survive \texttt{DISABLE TRIGGER ALL}.

\begin{table}[t]
\centering
\caption{Bypass scenarios and outcomes. From \texttt{scripts/bypass.sh}, \texttt{scripts/update\_forgery.sh}, \texttt{scripts/delete\_forgery.sh} transcripts and README threat-model enumeration.}
\label{tab:bypass}
\small
\begin{tabular}{p{4.7cm}cc}
\toprule
scenario & KNDB & trigger \\
\midrule
\texttt{DISABLE TRIGGER ALL}                & blocked & bypassed \\
\texttt{session\_replication\_role}         & blocked & bypassed* \\
\texttt{COPY} single-row (CIM\_SINGLE)      & blocked & blocked \\
\texttt{COPY} batch (CIM\_MULTI, post-F18)  & blocked & bypassed \\
\texttt{INSERT}, \texttt{INSERT SELECT}     & blocked & blocked \\
\texttt{UPDATE} of prefix (post-F20)        & blocked & bypassed \\
\texttt{DELETE} incumbent (post-F20)        & blocked & bypassed \\
\texttt{ON CONFLICT DO UPDATE} (post-F21)   & blocked$^{\dagger}$ & n/a \\
\texttt{ON CONFLICT DO NOTHING} (post-F21)  & blocked$^{\dagger}$ & n/a \\
Logical replication apply                   & blocked & bypassed* \\
\midrule
\multicolumn{3}{l}{Out of scope:} \\
\texttt{SET ACCESS METHOD heap}             & \multicolumn{2}{c}{schema-change threat} \\
\texttt{TRUNCATE}, \texttt{CLUSTER}, \texttt{VACUUM FULL} & \multicolumn{2}{c}{DDL-privileged; see Table~\ref{tab:tableam-audit}} \\
\bottomrule
\end{tabular}\\[2pt]
\noindent{\footnotesize
$^{\dagger}$ Blocked by the F21 override; but the SQL surface for this attack is not currently reachable because unique/exclusion constraint creation on epistemic tables fails at \texttt{heap\_getnext}'s \texttt{rd\_tableam} identity check (\texttt{heapam.c:1352}), so \texttt{ON CONFLICT (col)} has no arbiter to bind. Load-bearing is proven via the C-level probe \texttt{epistemic.\_probe\_speculative\_insert} and the source-rebuild disable-and-test in Section~\ref{sec:eval-f21}. The overrides are defensive: the engine-in-storage claim must not depend on that accidental index-support gap persisting.}
\end{table}

\subsection{Batch-size ceiling under COPY}
\texttt{scripts/bypass.sh} sweeps $N \in \{100, 1000, 10000, 20000\}$ via real \texttt{\textbackslash copy} at \texttt{max\_locks\_per\_transaction=64} against an epistemic table. $N=100$, $N=1000$, and $N=10000$ succeed; $N=20000$ raises \texttt{ERROR 53200 out of shared memory} at \texttt{LockAcquireExtended, lock.c:1080}. With the F18 \texttt{multi\_insert} override disabled (rebuild), the same $N=20000$ COPY succeeds, because \texttt{heap\_multi\_insert} takes no advisory locks; the ceiling is a direct consequence of routing through per-row enforcement, not an artefact of PostgreSQL machinery. Raising the GUC to 1024 pushes the ceiling to $\sim$250\,000 rows.

\subsection{Contention control (YCSB, F9 gate)}
Stage-1 gate results on YCSB-A at RC, 8 clients (\texttt{bench/results/summary/gate.csv}): KNDB epistemic 2057--2200 tps median (theta 0.0, 0.5), abort rate 0.050 and 0.092 with NEW\_LOSES the only abort family (no 40001 seen in the gate window). \texttt{pg\_heap} runs faster at 6600--7100 tps (no arbitration, no locks) but with no correctness guarantee. Stage-2 correctness cells (Section~\ref{sec:evaluation}) show that raw \texttt{pg\_heap} throughput comes with tens of thousands of duplicate live rows per 20-second window; correct-goodput (tps $\times$ (1 -- abort rate) $\times$ correctness) is 0 for \texttt{pg\_heap} on every Stage-2 cell.

\subsection{Temporal-shaped benchmarks (honest zero)}
KNDB scores 0.000 at $c=1$ on MemoryAgentBench Conflict\_Resolution, LongMemEval knowledge-update pairs, and MQuAKE-CF-3k edit chains (\texttt{bench/results/summary/stage3.md}). The reason is straightforward: these benchmarks expect last-writer-wins, and KNDB's F8 tiebreak is deliberately first-committer-wins. Both writes in each pair carry (MEASURED, spec=0, conf=1.0); the lattice cannot differentiate them by rank, so the tiebreak is what decides, and it is the wrong tiebreak for these workloads. This is disclosed in Section~\ref{sec:limitations}. On the same workloads \texttt{pg\_lww} scores 1.000 at $c=1$ and drops to 0.11--0.27 at $c=8$; the LWW ``win'' at $c=1$ is a determinism artefact of single-writer arrival order that vanishes under any contention.

%%% ---------------------------------------------------------------
\section{Related work}
\label{sec:related}
%%% ---------------------------------------------------------------

\textbf{Truth discovery.} The four TD algorithms in the evaluation cover the field's canonical designs. TruthFinder~\cite{yin2007truthfinder} is the seminal fixed-point on (source trust, fact confidence). CRH~\cite{li2014crh} formulates truth discovery as joint minimisation with a log-ratio weight update that is uniquely competitive on Zheng d\_sentiment below saturation, as our evaluation shows. CATD~\cite{li2015catd} adds a confidence-interval bound via $\chi^{2}_{\alpha/2, n}$ that is well-suited to long-tail source distributions. ACCU~\cite{dong2009accu} is Bayesian MAP over per-source accuracy and, in variants with copy detection, was originally designed to address adversarial data. We use base ACCU (no copy detection) and note in Section~\ref{sec:limitations} that copy-detection variants may shift the threshold. The Zheng et al.\ VLDB 2017 survey~\cite{zheng2017crowdsourcing} is the source of the d\_sentiment dataset and the reproducibility baseline we validated the reimplementations against. The Sybil-collapse framing draws on the same monotonicity structure the field has understood for a decade; the paper contribution is not the framing but the engine-level orthogonality that removes the collapse threshold entirely (Section~\ref{sec:formalism}).

\textbf{Provenance in PostgreSQL.} ProvSQL~\cite{senellart2018provsql} is the mature PostgreSQL extension for semiring provenance and probabilistic queries, tracked at row-set granularity via rewriting rather than in-storage. A recent PW25 demonstration~\cite{widiaatmaja2025provsql} adds update provenance through temporal databases with support for time-travel and undo. KNDB and ProvSQL do not overlap: ProvSQL propagates provenance through query evaluation; KNDB decides survivor selection at write time using a kind-primary total order. The two mechanisms are orthogonal and, in principle, composable.

\textbf{Agent-memory and epistemic warrant.} A recent line of position papers has argued that LLM agent memory is not, in practice, a database (in the sense of write-time correctness) and that the epistemic warrant conferred by an agent's tool boundary is thinner than commonly assumed. Orogat and Mansour~\cite{orogat2026agentmemory} argue for rethinking data foundations for long-term AI agent memory. Romanchuk and Bondar~\cite{romanchuk2026laundering} argue that tool boundaries alone do not confer epistemic warrant.

TOKI~\cite{wang2026toki} is our closest sibling work and warrants direct comparison. TOKI is a theory-first bitemporal operator algebra layered over an unmodified database engine: contradictions between competing memory assertions are resolved by projecting each assertion onto valid-time and system-time axes and picking the temporally-consistent survivor. TOKI has no epistemic-kind axis; every assertion is treated as a value carrying only temporal metadata. TOKI explicitly defers threat-model integration to future work. KNDB's contribution is precisely that deferred piece plus an engine-level realisation: an orthogonal kind axis that is assigned by the storage engine at write time and therefore cannot be forged by an untrustworthy writer, implemented as a PostgreSQL table access method rather than as an operator layered above one. TOKI's temporal semantics and KNDB's epistemic kind axis are compatible along different axes and, in principle, composable.

\textbf{Serialisable snapshot isolation.} KNDB's concurrency behaviour under \texttt{SERIALIZABLE} is standard PostgreSQL SSI~\cite{ports2012ssi}. The AM does not add a slot-level SIRead lock; predicate locking is inherited from \texttt{heapam}. Our early proof-of-concept had wrapped \texttt{PredicateLockTID} on a synthetic slot-hash TID, and F1 audited that wrapper and found it decorative (the wrapper's supposed contribution was already provided by heap's \texttt{PredicateLockRelation} plus \texttt{CheckForSerializableConflictIn}). The wrapper was removed.

\textbf{PostgreSQL as a research substrate.} We build on the PostgreSQL 18 table access method interface~\cite{postgres18tam}; the API surface has been stable across three major versions and is a productive vehicle for storage-layer research that must interoperate with a real query planner, executor, WAL, and recovery pipeline. Because rows on disk are plain heap tuples, existing PostgreSQL tools (\texttt{pg\_dump}, \texttt{pg\_upgrade}, logical replication, \texttt{pg\_basebackup}) work on KNDB tables without change.

%%% ---------------------------------------------------------------
\section{Limitations}
\label{sec:limitations}
%%% ---------------------------------------------------------------

We are explicit about where KNDB does not win.

\textbf{Truth discovery wins sub-saturation.} On Zheng d\_sentiment below the density-saturation cell (N < 20), CRH matches or beats KNDB on the confidence-forgery workload, and ACCU peaks above KNDB at N=5 (1.000 versus 0.927). CRH's log-ratio, CATD's $\chi^{2}$ confidence bound, and ACCU's Bayesian MAP each identify perfectly-wrong Sybils correctly when the coalition is small relative to per-slot honest votes. The paper's contribution is not that KNDB beats every TD baseline on every cell; it is that KNDB is the only system whose $k^{*}$ is unbounded, which becomes decisive at and above the saturation cell.

\textbf{Zero on temporal knowledge-editing benchmarks.} On MemoryAgentBench Conflict\_Resolution (37\,820 writes), LongMemEval knowledge-update (156 writes), and MQuAKE-CF-3k (12\,030 writes), KNDB scores 0.000 at $c=1$ because these benchmarks expect the later write to win and KNDB's tiebreak is first-committer-wins. A dedicated ``epistemic table with last-writer-wins tiebreak'' would be a straightforward configuration knob, but the current implementation does not expose it. Naive \texttt{pg\_lww} scores 1.000 on these benchmarks at $c=1$ and collapses to 0.11--0.27 at $c=8$, so LWW is not a general answer either; the tiebreak choice is a workload-specific configuration and no fixed choice is right for all workloads.

\textbf{Reputation farming defeats $k^{*}=\infty$.} The $k^{*}=\infty$ argument for KNDB against Sybil coordination depends on adversarial identities appearing fresh in the write stream, so that engine-assigned kind places them in \texttt{INFERRED} rather than \texttt{MEASURED}. A determined adversary who knows the tier-mapping rule can farm the reputation signal to reach Tier A before attacking: on Zheng, complete the qualification test (question ids 2000\ldots 2019) with correct answers so \texttt{quali\_acc} matches Tier A; on Book-Author, publish \texttt{n\_listings} volume on benign books to enter the top-K by source volume. Once inside Tier A, the adversary's writes are \texttt{MEASURED} and the lattice defends nothing. This is a real boundary of the mechanism. The formal argument holds for the write-time snapshot; long-lived reputation-farming requires an external audit or a stricter tier-assignment policy (e.g., cryptographically-signed source attestations, KYC on new bookstore registrations, or an audit that revokes Tier A on posthoc discovery of adversarial behaviour) that is out of scope for this paper. The two attack strategies that KNDB \emph{does} defend against are: (1) an adversary who has not farmed reputation and appears in the write stream in \texttt{INFERRED} regardless of self-reported confidence, which is the confidence-forgery workload of Section~\ref{sec:evaluation}, and (2) an adversary who has \texttt{INSERT} plus \texttt{ALTER TABLE} on the target table but no ability to change its tier assignment, which is the write-path bypass surface of Section~\ref{sec:motivation}.

\textbf{R2 and R5 SPI cost per non-MEASURED insert.} Both R2 (source resolution against \texttt{epistemic.source\_registry}) and R5 (slot-kind check against \texttt{epistemic.slot\_kind}) run a full \texttt{SPI\_connect} / \texttt{SPI\_execute\_with\_args} / \texttt{SPI\_finish} cycle per candidate row. Both run BEFORE the F6 per-slot advisory lock is acquired (\texttt{epistemic\_am.c} steps 1 vs 2), so neither introduces a deadlock risk against other advisory-lock holders. The per-row cost is measurable but small: a microbenchmark of 5\,000 INFERRED inserts versus 5\,000 MEASURED inserts on a fresh cluster puts R2's added overhead (one \texttt{count(*)} SPI call with a one-element text-array parameter) at $\sim$7 microseconds per row, or roughly 35\% on top of the $\sim$19 microsecond MEASURED baseline (which itself pays R5's SPI plus the overlap scan). Backend-local caching of \texttt{source\_registry} keyed off \texttt{CacheRegisterRelcacheCallback} would drive R2's per-row cost to near zero after a warm start and is a straightforward optimisation for a follow-up paper; we do not ship it here because the current cost is well below any workload's practical throughput floor and the caching adds an invalidation contract we would want to test more thoroughly than the F20 time budget allowed.

\textbf{Batch-size ceiling of $\sim$15\,000 rows per transaction at default GUC.} Documented in Section~\ref{sec:implementation}. Raising \texttt{max\_locks\_per\_transaction} to 1024 (a \texttt{postgresql.conf} change requiring restart) pushes the ceiling to $\sim$250\,000. Applies to COPY equally, because the F18 \texttt{multi\_insert} override routes batched COPY through the same per-row advisory lock. We chose not to drop the advisory lock in the batch path because doing so silently re-opens the F6 concurrency integrity leak.

\textbf{\texttt{ALTER TABLE SET ACCESS METHOD heap} out of scope.} The table owner can rewrite the relation onto plain heap, at which point the TAM callback is no longer bound to the relation and every subsequent write goes through unmodified heap. This is a schema-change threat, not a write-path threat. It is disclosed here and in the README threat-model section. Restricting the DDL is out of scope for this proof of concept; the natural next step is a PostgreSQL security policy that denies \texttt{SET ACCESS METHOD} away from \texttt{epistemic} for tables the epistemic contract is claimed over.

\textbf{Custom WAL rmgr is annotation only.} Recovery works byte-for-byte with the extension's WAL emitter disabled. Durability of the row comes from heap's \texttt{XLOG\_HEAP\_INSERT}. The custom rmgr 128 is retained as a named channel for a future logical decoding consumer to hook, but does not itself contribute to crash recovery. PostgreSQL's stock \texttt{pg\_waldump} does not load custom rmgrs and renders these records as \texttt{custom128 UNKNOWN}. This is disclosed.

\textbf{Eviction atomicity is PostgreSQL's.} The three-step eviction (winner insert, incumbent \texttt{sys\_time} close, audit row) executes inside one top-level PostgreSQL transaction. Atomicity is inherited, not novel. What the TAM contributes is that the three steps are inside the \texttt{tuple\_insert} callback and cannot be forgotten by a user-space writer, which is a convenience property, not a novel durability property.

\textbf{Seven of 42 TAM callbacks overridden.} We override \texttt{tuple\_insert}, \texttt{multi\_insert}, \texttt{tuple\_update} (F20), \texttt{tuple\_delete} (F20), \texttt{tuple\_insert\_speculative} (F21), \texttt{tuple\_complete\_speculative} (F21), and \texttt{relation\_toast\_am}. Every other callback (scan, index build, VACUUM, analyse, replication copy, cluster, tuple lock, ttids, and so on) is heap's, unmodified. Table~\ref{tab:tableam-audit} in Section~\ref{sec:tableam-audit} classifies all 42 callbacks with citable reasons. We inherit heap's behaviour on the 35 non-overridden ones, including its bugs. This is deliberate; the AM is a write-time enforcement point, not a storage-format replacement.

\textbf{Serialisable-level guarantees are PostgreSQL's.} Under \texttt{SERIALIZABLE}, KNDB inherits heap's predicate locking (\texttt{PredicateLockRelation}) rather than a slot-level SIRead lock. Two writers on non-overlapping slots may still hit an SSI abort because the predicate lock granulates to relation-level. This is a known SSI property in PostgreSQL~\cite{ports2012ssi}; adding slot-level predicate locking would require modifying \texttt{predicate.c} (a new lock target type) and is out of scope.

%%% ---------------------------------------------------------------
\section{Conclusion}
%%% ---------------------------------------------------------------

KNDB reports an engineering result: seven overrides in the PostgreSQL 18 table access method interface, an engine-assigned epistemic kind column, a per-slot advisory transaction lock, and a first-committer-wins tiebreak on true precedence ties are sufficient to make PostgreSQL survive a confidence-forgery attack that flattens every conventional PostgreSQL baseline. On two workloads with entirely different provenance shapes (source-tier data fusion, per-worker crowdsourcing qualification), the win is 63 points and 92.7 points over a confidence-only arbitrator and is proven load-bearing on the kind axis by source-rebuild disable-and-test. Against classic truth-discovery baselines the picture is more nuanced: KNDB is competitive below the per-dataset density-saturation cell and dominant at and above it, because kind rank is engine-assigned from independent metadata that adversarial identities do not possess. The threshold theorem $k^{*} \approx \rho_{\text{alg}} \cdot h_{\text{top}}$ makes the Sybil-collapse cell predictable from per-slot honest-support alone. What KNDB does not do is win on workloads whose ground truth is last-writer-wins; we are explicit about that, and about the batch-size ceiling, the schema-change threat, the annotation-only WAL role, and the reputation-farming boundary of the $k^{*}=\infty$ argument, throughout the paper. The value of this work, in our view, is the demonstration that an epistemic-kind primary sort key is deployable in a mainstream OLTP database in roughly 2400 lines of new C at the TAM layer, no core patch required, with a completeness argument over the full 42-callback TAM interface (Table~\ref{tab:tableam-audit}).

\begin{acks}
The authors used a large language model as a coding assistant and evaluation harness developer during this work. All system-level design, empirical results, formal arguments, and paper prose were reviewed and validated by the authors.
\end{acks}

\bibliographystyle{ACM-Reference-Format}
\bibliography{refs}

\end{document}
